\subsection{Model Validation}

Validating BNs is an essential step to verify whether the models, as a whole, adequately represent the domain of interest. Researchers can conduct this validation using different approaches, depending on data availability and the model's purpose. The methods identified in the \textbf{model validation} phase were organized into two main categories: expert-based validation and data-driven validation.

\vspace*{0.5\baselineskip}
\textbf{Methods for expert-based validation.} This type of validation verifies the coherence and plausibility of the model through the opinions of professionals with domain expertise. Table~\ref{tab:metodos_validacao_especialistas} lists the methods used in expert-based validation, which are described next.

\begin{table}[t]
\caption{Methods used for expert-based validation (Activity 3.1)}
\label{tab:metodos_validacao_especialistas}
\centering
\small
\setlength{\tabcolsep}{6pt}
\renewcommand{\arraystretch}{1.2}

\begin{tabularx}{0.75\linewidth}{>{\raggedright\arraybackslash}p{0.20\linewidth}
                              >{\raggedright\arraybackslash}X}
\toprule
\textbf{Method} & \textbf{Studies} \\
\midrule
Model walkthrough
& S1, S3, S4, S6, S7, S10, S18, S21, S28, S52, S73--S76, S81, S87, S93 \\
Sensitivity analysis
& S2, S18, S27, S37, S77 \\
Focus groups
& S74 \\
Case study
& S8, S36, S48, S80, S94, S103, S104 \\
Proof of concept
& S2, S27, S29, S33, S35, S43, S53, S58, S67, S88, S101 \\
\bottomrule
\end{tabularx}
\end{table}

\textit{Model walkthrough.} In studies such as S3, S4, and S6, domain experts created realistic scenarios to assess whether the BN's probability predictions aligned with the judgments they would make based on their own experience. For example, considering the effort estimation models presented in S4, domain experts examined several hypothetical scenarios and verified whether, for each scenario, the effort category with the highest probability predicted by the model matched the category they would have selected through subjective estimation. Some studies, such as S1 and S7, also referred to this approach as simulated scenarios.

\textit{Sensitivity analysis.} This method can be used to examine whether the model's behavior aligns with domain knowledge. In this process, experts analyze how changes in the states of selected nodes affect the probabilities of target variables and assess whether these changes follow expected patterns. When the model responds to such variations in ways that contradict expert intuition or experience, this may indicate problems in the BN's structure or parameters. Studies S18 and S27 applied sensitivity analysis to their models.

\textit{Focus groups.} The authors of S74 conducted a focus group with eight professionals from a company to evaluate the acceptance of the proposed solution based on the Technology Acceptance Model (TAM). During the session, the participants examined six simulated scenarios in which \added[id=A]{the BN nodes were instantiated to illustrate how predicted project outcomes varied across configurations of risk mitigation strategies}. The participants then completed a TAM-based questionnaire using a five-point Likert scale to assess two core dimensions: perceived usefulness and perceived ease of use. The activity lasted two and a half hours and combined collective discussion with structured elicitation of practitioners' perceptions.

\textit{Case study.} Other studies evaluated their models from the same perspective adopted in the focus group of S74, i.e., with an emphasis on usefulness and usability. For example, S8 evaluated its BN through a case study conducted with three agile software development teams over three sprints. The evaluation focused on assessing the practical usefulness of the proposed model for supporting the diagnosis and improvement of teamwork quality. The researchers collected input data for the model through interviews with the Scrum Masters, and they compared the results generated by the AgenaRisk tool with the teams' actual performance. They gathered participants' acceptance and perceptions of the model's applicability through Likert-scale questionnaires and open-ended questions, and analyzed these data using thematic analysis. The results indicated that the teams considered the model consistent with their reality, useful for identifying improvement opportunities, and easy to apply in the agile context, despite limitations in tool access.

The authors of S104 evaluated their BN model following the technology transfer model proposed by Gorschek et al.~\cite{gorschek2006model}, conducting three evaluation stages: an initial evaluation in an academic environment, a static validation with industry practitioners (both using TAM-based questionnaires), and a dynamic validation through a case study in a real industrial context. In this final stage, the authors applied the model throughout the life cycle of a software project to support Defect Causal Analysis (DCA) sessions based on real defect data. The evaluation considered both acceptance of the approach and the outcomes of the sessions. The participants expressed confidence in the model's inferences and identified systematic causes for the analyzed problems. 

\textit{Proof of concept.} Some studies used a proof of concept as a preliminary way to evaluate their BNs. This approach aims to demonstrate the feasibility and practical applicability of the proposed model through conceptual demonstrations or simulations. In general, a proof of concept shows the model in operation, illustrating how it can support decision-making, perform useful inferences, or respond to different situations, even in the absence of formal empirical validation. For example, the authors of S29 applied the BN to two decision scenarios to illustrate how different combinations of requirements and constraints (such as quality, functionality, effort, and time) influence the model's predictions, thereby demonstrating its usefulness for trade-off analysis. S88 used a proof of concept to apply its test case prioritization approach to eight consecutive versions of Apache Ant, exploring the model's behavior in a semi-automated environment and comparing multiple techniques. The practical demonstration of the BN served as initial evidence of its effectiveness and applicability, as the authors argued that the model would perform better when applied in practice to software systems whose failure rates are typically higher than in experimental settings.

\vspace*{0.5\baselineskip}

\textbf{Methods for data-driven validation.} Data-driven validation concerns with analyzing model performance by comparing the inferences generated by the BN with observational or historical data. This approach assesses the model's ability to reproduce real patterns, identify deviations, and quantify its accuracy. Table~\ref{tab:metodos_validacao_dados} lists the methods used for the data-driven validation of the reported models, which are described next.

\begin{table}[t]
\caption{Methods used for data-driven validation (Activity 3.2)}
\label{tab:metodos_validacao_dados}
\centering
\small
\setlength{\tabcolsep}{6pt}
\renewcommand{\arraystretch}{1.2}

\begin{tabularx}{\linewidth}{>{\raggedright\arraybackslash}p{0.30\linewidth}
                              >{\raggedright\arraybackslash}X}
\toprule
\textbf{Method} & \textbf{Studies} \\
\midrule
Predictive accuracy
& S3--S6, S8, S26, S27, S28, S40, S48, S59, S75, S78, S113, S114 \\
Cross-validation
& S12, S23, S24, S44, S68, S103, S107, S108, S116 \\
Statistical measures
& S12, S14, S20, S21, S23, S24, S27, S31, S52, S54, S59, S72, S77, S85, S92, S95, S96, S103, S107, S110, S111, S113, S118 \\
Comparison with existing models
& S38, S54, S77, S92, S95, S107, S108 \\
\bottomrule
\end{tabularx}
\end{table}

\textit{Predictive accuracy.} This validation method assesses the predictive capability of a BN model by comparing its predictions with real data from past projects. The model achieves accuracy when the category with the highest estimated probability for a target factor (e.g., effort) matches the actually observed value. In S4, the authors used predictive accuracy to validate effort estimation models. They input data from completed Web projects with known actual effort into the models, then compared their predictions --- the effort category with the highest probability --- with each project's actual effort. When discrepancies arose, they made adjustments and reprocessed the validation projects to ensure that the calibrations progressively improved the models. Other studies, such as S3 and S75, referred to this method as outcome adequacy.

\textit{Cross-validation.} Some studies used cross-validation to validate the constructed BN models. For example, the authors of S23 split the dataset into ten equally sized subsets (k-fold cross-validation). In the $i-th$ iteration $(i = 1, ..., k)$, the $i-th$ subset is used to test the BN learned from all remaining subsets. The evaluation considered two criteria: (i) accuracy rate, which indicates the ratio of correct predictions, and (ii) Fisher's exact test, used to verify whether a statistically significant correlation exists between the predicted and observed results. In S12, the authors validated the model using empirical data via k-fold cross-validation and summary statistics from the Weka tool.

\textit{Statistical measures.} To evaluate model performance, some studies used metrics computed by comparing actual and predicted results. For example, the authors of S12 used the Magnitude of Relative Error (MRE), Mean Magnitude of Relative Error (MMRE), Prediction at Level m (Pred.(m)), Accuracy, Mean Absolute Error (MAE), Root Mean Squared Error (RMSE), Relative Absolute Error (RAE), Root Relative Squared Error (RRSE), and Confusion Matrix.

\textit{Comparison with existing models.} Some studies validated their BN models using a data-driven approach that compares them with existing models, applying the same datasets and statistical measures to assess predictive accuracy. In S95, the authors compared the proposed model with others from the literature, using measures such as Mean Magnitude of Relative Error (MMRE), Balanced Mean Magnitude of Relative Error (BMMRE), Root Mean Square Error (RMSE), Normalized Root Mean Square Error (NRMSE), and the coefficient of determination $R^2$ on data from six software projects. In S108, the authors used 10-fold cross-validation to assess the model's predictive accuracy and compared it with classical data mining algorithms, including logistic regression, decision trees (C4.5), Naïve Bayes, TAN, and BAN, all implemented in Weka.

\vspace*{0.5\baselineskip}

\textbf{Comparative Analysis with the Methodologies Considered in the Construction of the Checklist.} One of the quality criteria adopted in the cross-domain review reported in the main article requires that the articles report the application of the proposed methodology. During the extraction of the essential activities from each BN methodology, we also identified the methods used to carry them out (see~\cite{rique2026supplementary}). Comparing these methods with those used in SE studies revealed both convergences and distinctive characteristics.

In the structure building phase, common methods include interviews (P1, P10, P13, P25), workshops and expert panels (P16, P25, P28, P29), surveys (P13, P23, P25), literature reviews (P7, P15, P22), expert reviews (P1, P2, P4, P7, P8), and the Delphi method (P17). Among the methods identified only in the BN methodologies are causal maps (P10, P22), Interpretive Structural Modeling (ISM) (P22), and Multi-Criteria Decision Analysis (MCDA) approaches\footnote{One SE study used AHP for structure building, but \added[id=A]{it was not included in the catalog because it did not meet the aggregate-score threshold.}} (P4, P15). In contrast, SE studies emphasize methods classified as reasoning mechanisms, such as the life cycle model, value model, quality model, and risk-centric modeling, which are absent in the analyzed BN methodologies and relate to domain-specific practices such as software life cycle management, value analysis, quality assessment, and risk management. Idioms (P1, P2) appear in both contexts, whereas the Parents and Children Search (PCS) algorithm (P8), used for structure learning, appeared exclusively in the analyzed methodologies. These results indicate that the structure building phase is the most dependent on the SE application domain.

In the parameter estimation and model validation phases, similarities prevail between the methods used in BN methodologies and those in SE studies. In the former phase, studies commonly use Expectation-Maximization (EM) (P1, P2, P7, P8, P29), the ranked nodes method (P4, P7, P15), noisy-OR and noisy-AND gates (P10), Maximum Likelihood Estimation (MLE) (P22), frequency quantification (P13, P23, P25, P28), and expert panels (P16). In the latter, sensitivity analysis (P1, P4, P13, P15, P16, P22, P25, P28, P29), scenario analysis (P10, P13, P15, P16, P17, P23, P25, P28), and cross-validation (P1, P2, P7, P8, P13) predominate.

Overall, many knowledge acquisition methods used in other domains also appear in SE. Domain dependence becomes more pronounced in the reasoning mechanisms used for structural development, which require adaptation to SE-specific models and practices, whereas activities related to parameter estimation and model validation rely on methods widely shared across fields.
